% Imporditud: tools/import_eksperimendid.py
% Allikas: Eesti füüsikaolümpiaadi eksperimendi ülesannete
%   kogu (Rael Kalda, 2023), ülesanne Ü59, lahendus L59.
\setAuthor{EFO žürii}
\setRound{lõppvoor}
\setYear{2012}
\setNumber{P E1}
\setDifficulty{2}
\setTopic{elektriahelad}
\setVahendid{patarei, kolm lampi pesades, juhtmed, ampermeeter ja voltmeeter}
\setPunktid{10}
\setAllikas{Eksperimendi kogu Ü59, lahendus lk 58}
\setLahendusseis{toores}

\prob{Lambid}
Koostage vooluring, milles kolmest ühesugusest hõõglambist kaks on ühendatud rööbiti ja kolmas lamp nendega jadamisi. Määrake, mitu korda on jadamisi ühendatud lambi takistus erinev ühes rööbiti ühendatud lambi takistusest. Millest on see erinevus tingitud?

\hint

\solu
Koostame allesitatud elektriskeemi.

Lülitame ampermeetri jadamisi lambiga $L_3$ ja voltmeetri rööbiti lambiga $L_3$. Mõõdame voolutugevuse lambis $L_3$ ja pinge selle klemmidel. Kuivõrd lambid on kõik ühesugused on voolutugevus lampides $L_1$ ja $L_2$ kaks korda väiksem kui lambis $L_3$. Me võime selle arvutada või mõõta paigutades ampermeetri skeemi jadamisi lambiga $L_1$. Paigutame voltmeetri rööbiti lambiga $L_1$ või $L_2$ ja mõõdame pinge lambi otstel. Kuivõrd lambid on ühesugused, piisab ühes mõõtmisest. Arvutame seosest R = U/I lampide takistused ja takistuste suhte. Lambi hõõgniidi takistus sõltub oluliselt hõõgniidi temperatuurist. Takistus suureneb temperatuuri suurenedes. Hõõgniidi temperatuur sõltub aga voolutugevusest lambis. Kuna rööbiti ühendatud lampides on voolutugevus väiksem kui nendega jadamisi ühendatud lambis, põlevad rööbiti ühendatud lambid tuhmimalt, nende hõõgniitide temperatuur on madalam kui lambi $L_3$ temperatuur ja nende taksitus on väiksem lambi $L_3$ takistusest.
\probend
